\documentclass[11pt]{article}

\usepackage{amsmath,amssymb,amsthm}
\usepackage[margin=1.1in]{geometry}
\usepackage{bm}
\usepackage{seqsplit}
\usepackage{array}
\usepackage[colorlinks=true,linkcolor=blue,citecolor=blue]{hyperref}

\newtheorem{lemma}{Lemma}
\newtheorem{proposition}{Proposition}
\newtheorem{remark}{Remark}
\newcommand{\R}{\mathbb{R}}
\newcommand{\eps}{\varepsilon}
\newcommand{\Bs}{B_{\mathrm{s}}}
\newcommand{\Bf}{B_{\mathrm{f}}}
\newcommand{\bnu}{\bm{\nu}}
\newcommand{\bN}{\mathbf{N}}
\newcommand{\bs}{\mathbf{s}}
\newcommand{\bv}{\mathbf{v}}
\newcommand{\bw}{\mathbf{w}}
\newcommand{\bx}{\mathbf{x}}
\newcommand{\jump}[1]{[\![#1]\!]}
\newcommand{\dS}{\,\mathrm{d}S}
\newcommand{\dV}{\,\mathrm{d}V}
\newcommand{\tr}{\operatorname{tr}}
\newcommand{\cof}{\operatorname{cof}}
\newcommand{\Div}{\operatorname{Div}}
\newcommand{\norm}[1]{\lVert #1 \rVert}

\title{The slipping fluid--solid interface in the referential\\
       formulation: kinematics, second variation, interface terms,\\
       discretisation and constraint enforcement}
\author{mfemElasticity method notes}
\date{}

\begin{document}
\maketitle
\sloppy

\noindent
This note is the reference for the slipping fluid--solid interface of
\texttt{\seqsplit{LinearQuasiStaticReferentialSelfGravitatingSlipProblem}}. The
general framework --- the referential description, the elastic tensor
dictionary, the gravity formulations and the volume operators of the
linearised referential system --- is
\texttt{doc/gravitating\_elasticity.md}; below, ``the referential
operators'' means the volume operators of its section ``The linearised
referential system''. The gauge penalty and its iterated refinement are
\texttt{doc/gauge\_penalty\_iteration.tex}; the welded alternative is
\texttt{doc/gauged\_fluid.md}; the per-class summary of the
implementation is \texttt{doc/quasi\_static\_models.tex}.

The question: when a fluid--solid interface slips, what does the
linearised static operator acquire on the interface, and how is the
slip discretised and its constraint enforced? Sections
\ref{sec:conv}--\ref{sec:collect} derive the second variation of the
total referential energy over the broken displacement space with the
linearised slip map, keeping every Jacobian factor explicit;
\S\ref{sec:disc} discretises the broken space; \S\ref{sec:enforce}
enforces the normal-jump constraint; \S\ref{sec:impl} and
\S\ref{sec:verif} describe the implementation and its verification;
Appendix~\ref{sec:verdict} compares the result with the published
treatments.

\emph{When these terms matter.} A tangential slip is gauge (removable
by relabelling the fluid) only if it extends into the fluid as a
relabelling field --- for a hydrostatic fluid, a field that is
divergence-free and tangent to the level surfaces of the equilibrium
potential, pressure and density. On an interface that is itself a level
surface this restricts the slip (rigid rotations of a spherical core
qualify, a general tangential slip need not); on an interface that is
not a level surface the construction generally fails. Tangential slip
is therefore not gauge in general (\texttt{doc/gauged\_fluid.md},
section ``Tangential slip and the welded space''). The slipping
formulation of this note is the general one: it allows a tangential
jump and enforces only normal continuity, and allowing a slip that is
not needed costs nothing in correctness. The welded (continuous-space)
gauged formulation is used where the slip it suppresses is absent or
removable, as for spherically symmetric models, where the two agree to
discretisation level (\S\ref{sec:verif}). Whenever the broken space is
used on a pressurised background the interface terms below are
obligatory: they vanish only where the equilibrium pressure does, as in
the gravity-free sliding tests.

\section{Conventions, and one warning}
\label{sec:conv}

Reference body $B = \Bs \cup \Bf$ with interface $\Sigma$; $\bN$ the
\emph{referential} unit normal on $\Sigma$, pointing out of the fluid;
$\dS$ the referential surface measure. Equilibrium mapping
$\varphi_e$ on the ball, $F_e = D\varphi_e$, $J_e = \det F_e$, and the
unnormalised \textbf{Nanson normal}
\begin{equation}
  \bnu := \cof(F_e)\,\bN = J_e F_e^{-T}\bN,
  \qquad
  \hat{\mathbf{n}}\,\mathrm{d}S_{\mathrm{phys}} = \bnu \dS,
  \qquad
  |\bnu|\dS = \mathrm{d}S_{\mathrm{phys}} .
  \label{eq:nanson}
\end{equation}

\begin{remark}[The Jacobian warning]\label{rem:jacobian}
By \eqref{eq:nanson} the two conventions ``unit physical normal
$\times$ physical measure'' and ``Nanson normal $\times$ referential
measure'' give the \emph{same} surface pairing. The error mode is
mixing them: $\hat{\mathbf{n}}$ against $\dS$ silently loses a factor
$|\bnu|$. Maitra \& Al-Attar (2024) drop exactly such a factor in
their fluid--solid boundary operator $Q$ --- an oversight, not a
choice of convention (harmless in their static applications, where the
term does not act). The working rule here: at this level of
complexity, \emph{explicit is better} --- every surface integral below
carries $(\bnu, \dS)$ with all factors written out, and no operator is
allowed to absorb a Jacobian tacitly.
\end{remark}

Motions $\varphi_{\mathrm{s}}$ on $\Bs$, $\varphi_{\mathrm{f}}$ on
$\Bf$. Stored energies: solid $W_{\mathrm{s}}(\bx, F)$; fluid
$W_{\mathrm{f}} = V(\bx, J)$ (bulk response only), with the referential
pressure
\begin{equation}
  \pi(\bx) := -\,D_J V(\bx, J_e)
  \qquad (= p^0 \circ \varphi_e \ \text{in Eulerian terms}),
\end{equation}
so the fluid first Piola--Kirchhoff stress at equilibrium is
$P_{\mathrm{f}} = -\pi\,\cof F_e$ and the fluid boundary traction per
referential area is $P_{\mathrm{f}}\bN = -\pi\,\bnu$: the
pressure--Nanson pairing appears by itself, Jacobian included.

Gravitational energy of the broken configuration, in the form that will
carry the whole derivation:
\begin{equation}
  E_g[\varphi_{\mathrm{s}}, \varphi_{\mathrm{f}}]
  = \operatorname{stat}_{\Phi}
    \Bigl[\, +\frac{1}{8\pi G}\int_{\R^d} |\nabla\Phi|^2 \dV
    \;+\; \sum_{r \in \{\mathrm{s,f}\}} \int_{B_r} \rho\,
          \Phi\circ\varphi_r \dV \Bigr],
  \label{eq:Eg}
\end{equation}
(the $+$ sign so that stationarity gives $\nabla^2\Phi = 4\pi G\varrho$
in our convention and the value is $\tfrac12\int\varrho\Phi$; the
sign is pinned by the gravity finite differences of \S\ref{sec:verif}).
A related point that motivates the whole construction: the
\emph{spatial} $\Phi$ is continuous ($C^1$) across the deformed
interface, but a referential potential composed \emph{region-wise with
each region's own motion} would not be --- at $\bx \in \Sigma$ the two
compositions evaluate $\Phi$ at the two distinct physical points
$\varphi_{\mathrm{s}}(\bx)$ and $\varphi_{\mathrm{f}}(\bx)$, so it
would jump by $\Phi(\varphi_{\mathrm{s}}) - \Phi(\varphi_{\mathrm{f}})
\approx \nabla\Phi_0\cdot(\text{slip})$ at linear order. Our $\zeta$
is single-valued \emph{by fiat}: it composes with the one globally
smooth extension $\tilde\varphi$, never with the fluid's own motion,
inheriting $\Phi$'s continuity (ordinary $H^1$ space, DtN untouched);
the price is the fluid-source mismatch below, which is the better
trade. All of these formulations --- Eulerian $\Phi$, referential
$a$-form with any smooth extension, single-valued or region-wise
$\zeta$ --- are \emph{exactly} equivalent, related by bijective changes
of variables in the stationary problem; the choices here are about
discretisation convenience only, and the extension-invariance is
asserted numerically at the solver level. Why start Eulerian rather than from the referential
$\tfrac{1}{8\pi G}\langle a(F)\nabla\zeta, \nabla\zeta\rangle$ form of
MA24? \emph{Not} because the $a$-form requires one composition map ---
it does not: the spatial field energy splits over the images of the
material regions (the attached broken motion tiles $\R^d$; the
attachment forbids gaps and overlaps), each piece transforms with its
\emph{own} motion, and the result is precisely the region-wise
$a(F_r)$-form in the broken variables $\zeta_r = \Phi\circ\varphi_r$
--- which \emph{is} the broken-$\zeta$ formulation, with the
continuity of $\Phi$, automatic Eulerian, resurfacing as the interface
compatibility $\zeta_{\mathrm{f}} = \zeta_{\mathrm{s}}\circ\sigma$ on
$\Sigma$ (composition on the side that $\sigma$ maps into, matching
$\varphi_{\mathrm{f}}|_\Sigma = \varphi_{\mathrm{s}}|_\Sigma\circ\sigma$;
linearised, this is the jump condition of \S\ref{sec:brokenzeta}). The reason to start Eulerian is derivational: keeping $\Phi$
spatial while varying makes its continuity and the gauge status of
the extension manifest, so the single-valued route's mismatch terms
and the broken route's jump condition are \emph{derived} rather than
posited; substituting $\Phi = \zeta\circ\tilde\varphi^{-1}$ afterwards
produces the $a(\tilde F)$-form \emph{plus} the fluid source composed
through $\tilde\varphi^{-1}\circ\varphi_{\mathrm{f}}$ --- which is
where the mismatch terms live. An $a$-form with a silently chosen
composition map would hide exactly this bookkeeping. Note that
in \eqref{eq:Eg} the Eulerian potential $\Phi$ is composed with
\emph{each region's own motion}, and its first variation is a pure
volume integral in each region
($\delta E_g = \sum_r \int \rho\, (\nabla\Phi\circ\varphi_r)\cdot
\bm{\eta}_r$): \textbf{gravity contributes no interface flux}, a fact
used repeatedly.

\section{Kinematics of the slip}
\label{sec:kin}

The fluid stays attached but slips:
$\varphi_{\mathrm{f}}(\Sigma) = \varphi_{\mathrm{s}}(\Sigma)$ as
surfaces. Parameterise
$\varphi_{\mathrm{f}}|_\Sigma = \varphi_{\mathrm{s}}|_\Sigma \circ
\sigma$ with a slip map $\sigma : \Sigma \to \Sigma$,
\begin{equation}
  \sigma_\eps = \mathrm{id} + \eps\,\bs + \tfrac{\eps^2}{2}\,\bm{\sigma}_2
  + O(\eps^3),
  \qquad \bs \in T\Sigma .
\end{equation}
The first-order coefficient is tangential, but the second-order one is
\emph{not} free to be: a family of maps that stays on a curved surface
has the forced normal component
\begin{equation}
  \bN\cdot\bm{\sigma}_2 = -\,\mathrm{II}(\bs, \bs),
  \qquad
  \bm{\sigma}_2 = \bs_2 - \mathrm{II}(\bs,\bs)\,\bN,
  \quad \bs_2 \in T\Sigma \ \text{free},
  \label{eq:sigma2}
\end{equation}
with $\mathrm{II}$ the referential second fundamental form of $\Sigma$
(circle of radius $r$, outward $\bN$: $\bN\cdot\bm{\sigma}_2 =
-|\bs|^2/r$, as an exact rotation family shows). Taking
$\bm{\sigma}_2$ tangential would lose the curvature term below; the
finite-$\eps$ verification families of \S\ref{sec:verif} resolve it.
Expand $\varphi_{\mathrm{s},\eps} = \varphi_e + \eps\,\bv_{\mathrm{s}}$
(a straight path; the second-order freedom sits in $\sigma$ and in the
fluid's interior). Composing and Taylor-expanding $\varphi_e$ and
$\bv_{\mathrm{s}}$ along the surface,
\begin{equation}
  \varphi_{\mathrm{f},\eps}\big|_\Sigma
  = \varphi_e
  + \eps\,\bigl(\bv_{\mathrm{s}} + F_e \bs\bigr)
  + \frac{\eps^2}{2}\,\underbrace{\bigl(
      2\,D\bv_{\mathrm{s}}[\bs] + D^2\varphi_e[\bs,\bs] + F_e \bs_2
    \bigr)}_{=:\ \mathbf{b}_2} + O(\eps^3).
  \label{eq:expand}
\end{equation}

\textbf{First order.} The fluid trace is
$\bv_{\mathrm{f}}|_\Sigma = \bv_{\mathrm{s}} + F_e\bs$, i.e.\ the jump
$\jump{\bv} := \bv_{\mathrm{s}} - \bv_{\mathrm{f}} = -F_e\bs$ lies in
$F_e(T\Sigma)$. Since $\bnu\cdot F_e\mathbf{t} = J_e\,\bN\cdot\mathbf{t}
= 0$ for $\mathbf{t}\in T\Sigma$, this is exactly
\begin{equation}
  \boxed{\ \bnu\cdot\jump{\bv} = 0 \ }
  \qquad\text{with}\qquad
  \bs = -F_e^{-1}\jump{\bv} .
  \label{eq:constraint}
\end{equation}
The normal-jump constraint of the sliding interface, with the
\emph{mapped} normal: at $\varphi_e \neq \mathrm{id}$ the constraint
normal is $\bnu$ --- which is what \texttt{Diffeomorphism::MapNormal}
and the mapped \texttt{BoundaryNormalNormalIntegrator} supply.

\textbf{Second order.} The datum forced on the fluid's second-order
boundary displacement is $\mathbf{b}_2$ of \eqref{eq:expand} with
$\bs_2$ replaced by $\bm{\sigma}_2$. Its $\bnu$-component: the free
tangential part $\bs_2$ drops ($\bnu \perp F_e T\Sigma$), while the
forced normal part contributes $\bnu\cdot F_e\bN\cdot(-\mathrm{II}) =
-J_e\,\mathrm{II}(\bs,\bs)$ (using $\bnu\cdot F_e\bN = J_e$), so
\begin{equation}
  \bnu\cdot\mathbf{b}_2
  = 2\,\bnu\cdot D\bv_{\mathrm{s}}[\bs]
  + \bnu\cdot D^2\varphi_e[\bs,\bs]
  - J_e\,\mathrm{II}(\bs,\bs).
  \label{eq:b2n}
\end{equation}

\begin{remark}[The curvature enters, and immediately re-hides]
\label{rem:wein}
For a tangent field $\bs$ on $\Sigma$, differentiating
$\bN\cdot\bs = 0$ gives $\bN\cdot D\bs[\bs] = -\mathrm{II}(\bs,\bs)$
--- Valette's (1986, 1991) Weingarten-operator content. Combined with
$D(F_e\bs)[\bs] = D^2\varphi_e[\bs,\bs] + F_e D\bs[\bs]$, the last two
terms of \eqref{eq:b2n} collapse:
\begin{equation}
  \bnu\cdot D^2\varphi_e[\bs,\bs] - J_e\,\mathrm{II}(\bs,\bs)
  = \bnu\cdot D(F_e\bs)[\bs]
  = -\,\bnu\cdot \nabla_\Sigma\jump{\bv}\,[\bs],
  \label{eq:collapse}
\end{equation}
since $F_e\bs = -\jump{\bv}$ and only tangential derivatives act.
Hence
\begin{equation}
  \bnu\cdot\mathbf{b}_2
  = 2\,\bnu\cdot \nabla_\Sigma\bv_{\mathrm{s}}[\bs]
  - \bnu\cdot \nabla_\Sigma\jump{\bv}[\bs]
  = \bnu\cdot \nabla_\Sigma(\bv_{\mathrm{s}} + \bv_{\mathrm{f}})[\bs]:
  \label{eq:b2clean}
\end{equation}
tangential derivatives of the two \emph{traces} only. Neither the
curvature nor $D^2\varphi_e$ is ever assembled --- the referential
analogue of the volume terms needing no $\nabla\nabla\Phi_0$ --- and
the interface integrator needs only surface shape-function derivatives
and $\bnu$.
\end{remark}

\section{First variation: equilibrium and the multiplier}
\label{sec:first}

For an unconstrained broken variation $\bm{\eta} = (\bm{\eta}_{\mathrm{s}},
\bm{\eta}_{\mathrm{f}})$, with outward normals $-\bN$ for $\Bs$ and
$+\bN$ for $\Bf$ on $\Sigma$,
\begin{equation}
  \delta E[\bm{\eta}]
  = \sum_r \int_{B_r} (\text{EL}_r)\cdot\bm{\eta}_r \dV
  + \oint_\Sigma \bigl[ (P_{\mathrm{f}}\bN)\cdot\bm{\eta}_{\mathrm{f}}
                      - (P_{\mathrm{s}}\bN)\cdot\bm{\eta}_{\mathrm{s}}
                 \bigr] \dS + (\text{outer terms}),
  \label{eq:firstvar}
\end{equation}
where $\text{EL}_r$ is the region's full (elastic + gravitational)
Euler--Lagrange density. Inserting the constrained variations
$\bm{\eta}_{\mathrm{f}} = \bm{\eta}_{\mathrm{s}} + F_e\mathbf{t}$
($\mathbf{t} \in T\Sigma$ arbitrary) and demanding stationarity:
\begin{align}
  \text{EL}_r &= 0 \quad\text{in } B_r
  &&\text{(interior equilibrium)},\\
  P_{\mathrm{s}}\bN &= P_{\mathrm{f}}\bN = -\pi\,\bnu \quad\text{on }
  \Sigma
  &&\text{(traction continuity; purely normal)}.
  \label{eq:tract}
\end{align}
The second line is the \emph{equilibrium consistency condition} on any
generated background: the solid side must hand the interface a normal
traction of density $-\pi\bnu$ per referential area. Writing the
constraint \eqref{eq:constraint} with a multiplier,
$\oint \lambda\, \bnu\cdot\jump{\bv} \dS$, the same computation
identifies
\begin{equation}
  \boxed{\ \lambda = \pi \ }
\end{equation}
(the referential pressure; by Remark~\ref{rem:jacobian} the
per-physical-area convention gives the same $\lambda$). This is the
multiplier of the \emph{nonlinear} (first-variation) problem: it states
the equilibrium traction condition \eqref{eq:tract}, which can be
checked directly on a generated background. The \emph{linearised}
problem solved by the library imposes the linearised constraint
\eqref{eq:constraint} on the perturbation; its multiplier is the
first-order interface normal traction $\varpi^{(1)}$ paired with
$\bnu$ (the traction \emph{perturbation}), not $\pi$. That is what the
augmented-Lagrangian multiplier and the KKT multiplier of
\S\ref{sec:enforce} converge to; $\pi$ itself enters the linearised
operator only through the background coefficients of $B_\Sigma$ and the
fluid Hessian below.

\section{The second variation}
\label{sec:second}

Differentiate $E(\eps) := E[\varphi_{\mathrm{s},\eps},
\varphi_{\mathrm{f},\eps}]$ twice along the constrained family of
\S\ref{sec:kin}:
\begin{equation}
  \frac{d^2 E}{d\eps^2}\Big|_0
  = \delta^2 E[(\bv_{\mathrm{s}}, \bv_{\mathrm{f}}),
               (\bv_{\mathrm{s}}, \bv_{\mathrm{f}})]
  \;+\; \delta E[(\mathbf{0}, \mathbf{u}_2)],
  \label{eq:d2E}
\end{equation}
where $\mathbf{u}_2$ is the fluid's second-order displacement (interior
values free, trace $\mathbf{b}_2$ forced). By \eqref{eq:firstvar} with
the equilibrium conditions, the second term collapses to a single
surface integral --- the interior part dies on $\text{EL} = 0$
(gravity's contribution included), the free interior values of
$\mathbf{u}_2$ therefore drop out (as they must: they are
parameterisation gauge), and
\begin{equation}
  \delta E[(\mathbf{0},\mathbf{u}_2)]
  = \oint_\Sigma (P_{\mathrm{f}}\bN)\cdot\mathbf{b}_2 \dS
  = -\oint_\Sigma \pi\; \bnu\cdot\mathbf{b}_2 \dS .
\end{equation}
With \eqref{eq:b2n} and $\bs = -F_e^{-1}\jump{\bv}$:

\begin{proposition}[The interface form]\label{prop:BSigma}
The entire interface contribution to the second variation is, by
\eqref{eq:b2clean},
\begin{equation}
  B_\Sigma(\bv, \bv)
  = -\oint_\Sigma \pi\;
      \bnu\cdot \nabla_\Sigma(\bv_{\mathrm{s}} + \bv_{\mathrm{f}})[\bs]
    \dS,
  \qquad \bs = -F_e^{-1}\jump{\bv}
  \label{eq:BSigma}
\end{equation}
--- the mean-trace derivative against the jump --- used in its
symmetrised polarisation
\begin{equation}
  B_\Sigma(\bv, \bv')
  = -\tfrac12\oint_\Sigma \pi \Bigl[
      \bnu\cdot \nabla_\Sigma(\bv_{\mathrm{s}} + \bv_{\mathrm{f}})[\bs']
    + \bnu\cdot \nabla_\Sigma(\bv_{\mathrm{s}}' + \bv_{\mathrm{f}}')[\bs]
    \Bigr] \dS .
  \label{eq:BSigmaPol}
\end{equation}
There are \emph{no other} elastic interface terms: the volume Hessians
$\delta^2 E$ are straight-path second derivatives of volume integrals
and generate none.
\end{proposition}

In particular there is no separate ``flux-variation $\times$ jump''
term: such a term arises only in a non-minimal organisation of the
same computation. Every coefficient of \eqref{eq:BSigma} --- the
curvature content of \eqref{eq:b2n} included --- is checked against
exact finite-$\eps$ slip families (\S\ref{sec:verif}).

\subsection{The volume Hessians, and the fluid dictionary re-derived}

The solid part of $\delta^2 E$ is the referential operator of
\texttt{doc/gravitating\_elasticity.md}, with the solid's own field. The fluid part, from $V(\bx, J)$ with
$H := F_e^{-1} D\bv_{\mathrm{f}}$, $\delta J = J_e \tr H$, and
$\delta^2 J = J_e[(\tr H)^2 - \tr(H^2)]$ along straight paths:
\begin{equation}
  \delta^2 E_{\mathrm{f}}
  = \int_{\Bf} \Bigl[ V_{JJ} J_e^2 (\tr H)^2
    \;-\; \pi J_e \bigl( (\tr H)^2 - \tr(H^2) \bigr) \Bigr] \dV .
  \label{eq:fluidhess}
\end{equation}

\begin{lemma}[$\mu_b = \pi$ from $V(\bx,J)$ alone]\label{lem:mub}
Expression \eqref{eq:fluidhess} equals the general material +
geometric split with the \emph{bare} fluid moduli: using
$\tr(H^2) = 2|\mathrm{sym}\,H|^2 - |H|^2$,
\[
  \eqref{eq:fluidhess}
  = \int_{\Bf} \Bigl[ \underbrace{\kappa_b (\tr H)^2
      + 2\pi\,|\mathrm{sym}\,H|^2}_{\text{material, } \mu_b = \pi}
      \;\underbrace{-\;\pi\,|H|^2}_{\text{geometric, } S_e = -\pi
      \text{-Piola}} \Bigr] J_e\text{-weighted} \dV
\]
with $\kappa_b$ the corresponding bulk combination of $V_{JJ}$ and
$\pi$. The pressure-induced shear rigidity $\mu_b = \pi$ of the elastic
tensor dictionary (\texttt{doc/gravitating\_elasticity.md}, section
``The elastic tensor dictionary'') is thus forced by the fluid's
constitutive class alone --- independently confirming the conversion
that the relabelling-invariance test
(\texttt{TestReferentialProblem}, \texttt{FluidRelabellingNullPair})
verifies numerically.
\end{lemma}

\subsection{\texorpdfstring{Gravity: single-valued $\zeta$, and two vanishing results}{Gravity: single-valued zeta, and two vanishing results}}
\label{sec:grav}

Let $\tilde\varphi$ be the solid motion extended smoothly over $\Bf$
(and the buffer), $\zeta := \Phi\circ\tilde\varphi$ single-valued on
the ball, and $\bw := \bv_{\mathrm{f}} - \tilde{\bv}$ the
fluid-vs-extension mismatch ($\bw|_\Sigma = -\jump{\bv} = F_e\bs$, so
$\bnu\cdot\bw|_\Sigma = 0$ by \eqref{eq:constraint}). Substituting
$\Phi = \zeta\circ\tilde\varphi^{-1}$ in \eqref{eq:Eg}: the field term
becomes the mapped $a(\tilde F)$-form of the referential operators;
the solid source is
$\int_{\Bs}\rho\,\zeta$; the fluid source composes through
$\psi_\eps := \tilde\varphi_\eps^{-1}\circ\varphi_{\mathrm{f},\eps}
= \mathrm{id} + \eps F_e^{-1}\bw + O(\eps^2)$:
\begin{equation}
  \int_{\Bf} \rho\; \zeta\circ\psi_\eps \dV
  = \int_{\Bf} \rho\,\zeta
  + \eps \int_{\Bf} \rho\,(F_e^{-1}\bw)\cdot\nabla\zeta
  + O(\eps^2).
  \label{eq:source}
\end{equation}

\begin{lemma}[First order: no gravity interface term]\label{lem:noflux}
Integrating the $O(\eps)$ term of \eqref{eq:source} by parts, the
boundary contribution is
$\oint_\Sigma \rho_{\mathrm{f}}\,(\bw\cdot\bnu/J_e)\,\zeta \dS = 0$
by the constraint. The welded result --- no gravity interface terms
--- survives the slip at first order; the volume remainder
$-\int \Div(\rho F_e^{-1}\bw)\,\zeta$ is the fluid relabelling source
(gauge for admissible $\bw$, physical otherwise), exactly as in the
gauged-fluid analysis.
\end{lemma}

\begin{lemma}[Second order: gravity stays volume-only]\label{lem:vol}
At $O(\eps^2)$ the composed source generates (i) terms linear in the
second-order data, which are precisely the $\delta E[\ddot\varphi]$
contributions already accounted in \eqref{eq:d2E} --- counting them
again would double-book; and (ii) the quadratic term
$\tfrac12\int\rho\, (F_e^{-1}\bw)\cdot \nabla\nabla\zeta\,
(F_e^{-1}\bw)$, whose second derivative of $\zeta$ is removed by the
standard by-parts, with boundary term proportional to
$\bnu\cdot\bw|_\Sigma = 0$ again. Hence the slip's gravitational
physics enters \emph{only} through volume terms on $\Bf$: the mismatch
coupling $\int_{\Bf}\rho\,(F_e^{-1}\bw)\cdot\nabla\zeta^1$ and the
$a'/a''$-machinery of the referential operators evaluated with the
extended field --- no
gravity interface integrals at any order kept.
\end{lemma}

There is no gravity interface coupling in the single-valued
organisation. The extension $\tilde{\bv}$ is gauge; the assembled operator's
independence of it (given the constraint) is the built-in invariance
test, exactly as for the vacuum extension.

\subsection{The gravity Hessian of the broken motion, explicitly}
\label{sec:gravhess}

For implementation the volume-only content of Lemma~\ref{lem:vol} must
be written out. Expanding $\psi_\eps = \tilde\varphi_\eps^{-1}\circ
\varphi_{\mathrm{f},\eps} = \mathrm{id} + \eps\,\mathbf{g}_1 +
\tfrac{\eps^2}{2}\mathbf{g}_2$ by composing the two motions
($\tilde\varphi = \varphi_e + \eps\tilde{\bv}$ straight):
\begin{equation}
  \mathbf{g}_1 = F_e^{-1}\bw, \qquad
  \mathbf{g}_2 = F_e^{-1}\bigl( \mathbf{u}_2
    - D^2\varphi_e[\mathbf{g}_1, \mathbf{g}_1]
    - 2\,D\tilde{\bv}[\mathbf{g}_1] \bigr),
  \qquad \bw := \bv_{\mathrm{f}} - \tilde{\bv} .
\end{equation}
The $\mathbf{u}_2$ piece of $\mathbf{g}_2$ is the already-booked
path-curvature contribution of \eqref{eq:d2E} (interior EL $\times$
$\mathbf{u}_2$ plus the boundary term that became $B_\Sigma$) and must
\emph{not} be counted again; the genuinely quadratic remainder,
together with the $\nabla\nabla\zeta^0$ term of the source expansion,
gives the extra gravity Hessian on $\Bf$. Killing
$\nabla\nabla\zeta^0$ by parts (boundary term $\propto
\mathbf{g}_1\cdot\bN = \bw\cdot\bnu/J_e = 0$) produces a
$\nabla_{\mathbf{g}_1}\mathbf{g}_1$ remainder that \textbf{cancels the
$D^2\varphi_e$ term identically} (via $\nabla_{\mathbf{g}_1}(F_e
\mathbf{g}_1) = D^2\varphi_e[\mathbf{g}_1,\mathbf{g}_1] +
F_e\nabla_{\mathbf{g}_1}\mathbf{g}_1$) --- the volume analogue of the
curvature collapse \eqref{eq:collapse}. With $\mathbf{b} :=
F_e^{-T}\nabla\zeta^0$, the additions to the referential operators
(which act
with the \emph{extended} field $\tilde{\bv}$ ball-wide) are, in
Hessian normalisation, the mismatch coupling $H_{\zeta\bw}$ and the
extra displacement Hessian $H_{\bw\bw}$,
\begin{align}
  H_{\zeta\bw}
  &= 2\int_{\Bf} \rho\,(F_e^{-1}\bw)\cdot\nabla\zeta^1 \dV ,
  \label{eq:mismatch}\\
  H_{\bw\bw}
  &= -\int_{\Bf} \Div(\rho F_e^{-1}\bw)\,
      \bigl((F_e^{-1}\bw)\cdot\nabla\zeta^0\bigr) \dV
   \;-\;\int_{\Bf} \rho\,\mathbf{b}\cdot D\bw[F_e^{-1}\bw] \dV
   \nonumber\\
  &\qquad\;-\;2\int_{\Bf} \rho\,\mathbf{b}\cdot
      D\tilde{\bv}[F_e^{-1}\bw] \dV,
  \label{eq:wextra}
\end{align}
all first-derivative volume forms (in the discrete realisation the
$\Div(\rho\,\cdot)$ form may be traded against a gradient of the
projected $\nabla\zeta^0$ field to avoid density gradients). The
$\tilde{\bv}$-dependence of \eqref{eq:wextra} and of the referential
operators
is extension gauge: the assembled operator is invariant under changes
of the extension for constrained fields --- the built-in invariance
test, as for the vacuum extension.

When $\zeta^1$ is eliminated at fixed motion, its stationary
contribution to the Hessian is \textbf{negative definite}
($-\tfrac{1}{4\pi G}\int|\nabla\zeta^{1*}|^2$): the induced potential
perturbation always lowers the energy --- self-gravity destabilises,
as it must.

\textbf{Verification} (\texttt{TestSlipInterface},
\texttt{GravityHessianIdentity}, at $\le 3\times10^{-4}$ relative):
with a rigid,
stress-free solid ($\rho_{\mathrm{s}} = 0$, $\pi(r_c) = 0$ hydrostatic
fluid, so $B_\Sigma$ and the $\tilde{\bv}$-terms switch off and
\eqref{eq:mismatch}--\eqref{eq:wextra} are isolated), exact
disc-preserving families give the total energy exactly --- the
gravitational part in closed 1-D form by the shell theorem, since the
families' densities stay radial. Verified: the radial-squeeze identity
(which pins the sign of the eliminated-$\zeta^1$ part: the opposite
sign misses the exact energy by twice the stationary value); a $\theta$-dependent azimuthal family (not
volume-preserving, hence \emph{not} a relabelling --- the assembled
prediction matched the exact energy to $4\times10^{-6}$ anyway); the
\emph{axisymmetric} rotation, which is the true relabelling and is
jointly annihilated with gravity on (elastic $\pi(r)$-terms,
$\nabla\nabla\zeta^0$-term and stationary part cancelling); the mixed
family; and the vanishing of the gravity--slip cross terms demanded by
Lemma~\ref{lem:vol}.

\textbf{Discrete realisation} (same file,
\texttt{DiscreteGravityPiecesMatchQuadrature}, against the same analytic
quadrature; order~1 at $\sim\!4\times10^{-2}$, order~2 at
$2$--$5\times10^{-4}$ relative on every piece). The extension is a
prescribed sparse operator $\tilde{\bv} = E\,\bv_{\mathrm{s}}$
(\texttt{NewRadialFluidExtension}): interface trace rows copy the
solid values exactly through the SubMesh dof pairing --- to round-off,
the property the mismatch $\bw = \bv_{\mathrm{f}} - \tilde{\bv}$ and
the vanishing lemmas rely on --- and interior rows take
$t(r)\,\bv_{\mathrm{s}}(r_c\hat{\bx})$ with $t = (r/r_c)^2$ vanishing
at the centre, where the radial direction is undefined; the interior
rule is gauge. The volume pieces at $\varphi_e = \mathrm{id}$ assemble
from the library integrators: the
$\int\rho\,\bw\cdot\nabla\nabla\zeta^0\,\bw'$ term as a vector mass
matrix $M$ with matrix coefficient $\rho\nabla\nabla\zeta^0$
($= 2\pi G\rho^2\,\mathbf{1}$ inside a uniform disc); the mismatch
coupling $2\int\rho\,\bw\cdot\nabla\zeta^1$ as the mixed
vector--scalar-gradient form with coefficient $\rho$; and the
$\tilde{\bv}$-term $-2\int\rho\,\nabla\zeta^0\cdot D\tilde{\bv}[\bw]$
as the \emph{compensated combination} $G - M$: the library's
vector--gradient--vector form interpolates the scalar
$\nabla\zeta^0\cdot\tilde{\bv}$ nodally and differentiates the
\emph{product}, so $G$ realises
$\int\rho\,\bw\cdot\nabla(\nabla\zeta^0\cdot\tilde{\bv})$ and the
excess $\int\rho\,\bw^T(\nabla\nabla\zeta^0)\,\tilde{\bv}$ is exactly
the mass term already assembled. The uncompensated form is wrong by
a quantity of the size of the target (on the test fields it cancels the
target almost entirely), so the compensation is essential, and the
discrete-versus-quadrature cross-check is what guards such
library-form semantics.

\begin{remark}[Topology, the extension, and the broken-$\zeta$
alternative]
With a solid inner core the fluid shell has two solid boundaries and
the extension interpolates between two \emph{independent} solid
motions (a spherical ICB/CMB pair even leaves the inner core's
rotation free --- known null-space bookkeeping). Nothing in the theory
changes: the constraint gives $\bnu\cdot\bw = 0$ on \emph{each
component} of $\partial\Bf$, the vanishing lemmas hold
interface-by-interface, and $B_\Sigma$ appears once per interface;
only the extension rule must interpolate two-sided data. Oceans are the
degenerate case: constraint at the seafloor only. The implemented
single-valued organisation supports one fluid region inside a solid
shell, with the one-sided radial extension of \S\ref{sec:impl}; for
aspherical fluid regions a prescribed geometric extension has no
natural recipe. A \emph{broken}-$\zeta$ formulation (region-wise
composition; jump condition $\jump{\zeta^1} =
-\bs\cdot\nabla_\Sigma\zeta^0$ per interface, and no fluid extension at
all) is exactly equivalent and scales trivially with topology: it is
the self-benchmark of the mismatch terms and the organisation used for
mapped and aspherical backgrounds. It is derived in full in
\S\ref{sec:brokenzeta}.

A third organisation, by superposition of \emph{sources}:
$\Phi = \sum_r \Phi_r$ with $\Phi_r$ generated by region $r$'s density
alone, so the energy splits into the pairwise interactions
$\tfrac12\sum_{r,r'}\int\rho_r\,\Phi_{r'}$ and each part can be
computed separately --- at evaluation level this is used by the
verification (its shell-theorem references are exactly such interaction
integrals). As a \emph{formulation} its cost is one
ball-wide potential block (with its own DtN) per region, and the extra
solves do not remove the interface structure: what breaks $\zeta$ is
not any discontinuity of the field --- each $\Phi_r$ is globally $C^1$
--- but the \emph{composition with the broken motion}, and each split
potential faces the same choice as the total one: composed region-wise
it jumps on every interface (with the partial field's gradient,
$\jump{\zeta_r^1} = -\bs\cdot\nabla_\Sigma\zeta_r^0$, summing back to
the total condition), composed through a single smooth extension it
generates mismatch terms per region pair. Its natural use is therefore
diagnostic --- region-wise attribution of the induced potential ---
not primary.
\end{remark}

\begin{remark}[Stratification]
For a two-parameter fluid the energy carries the additional advected
scalar; its second variation contributes the interface-adjacent
$N^2$-type stability term through the \emph{volume} density
$\rho(\bx)$-class terms restricted by the smaller relabelling class
--- again no new surface integral, but a genuine change to which
$\bw$'s are gauge. That bookkeeping belongs to a two-parameter
constitutive model, not to the interface; the library's rheologies carry
no second advected parameter.
\end{remark}

\subsection{\texorpdfstring{Gravity: the broken-$\zeta$ organisation}{Gravity: the broken-zeta organisation}}
\label{sec:brokenzeta}

The alternative anticipated above, in full. Compose the potential
region-wise with each region's \emph{own} motion, $\zeta_r :=
\Phi\circ\varphi_r$, where $r$ runs over the smoothness components of
the broken motion: each fluid region with its own motion, and each
solid component together with its vacuum continuation where it abuts
the buffer (the outer region's composition map must be smooth across
the free surface and equal the identity at the DtN sphere, so the
\emph{vacuum} extension of the welded problem survives, gauge as
always; the \emph{fluid} extension does not). Splitting the Eulerian
field energy of \eqref{eq:Eg} over the deformed regions --- a tiling:
the attachment forbids gaps and overlaps --- and transforming each
piece with its own motion,
\begin{equation}
  E_g = \operatorname{stat}
  \Bigl[\, \sum_r \frac{1}{8\pi G}\int_{B_r}
  \langle a(F_r)\nabla\zeta_r, \nabla\zeta_r\rangle \dV
  \;+\; \sum_r \int_{B_r}\rho\,\zeta_r \dV \Bigr],
  \label{eq:Ebroken}
\end{equation}
stationary over the broken $\zeta$ subject to the compatibility
$\zeta_{\mathrm{f}} = \zeta_{\mathrm{s}}\circ\sigma$ on $\Sigma$ (the
continuity of $\Phi$ at the deformed interface, \S\ref{sec:conv}).
Both sources in \eqref{eq:Ebroken} are \emph{exact}: no composition
is ever expanded, so the entire mismatch machinery of
\S\ref{sec:gravhess} --- the fluid extension, $\bw$, and
\eqref{eq:mismatch}--\eqref{eq:wextra} --- disappears identically.
Its place is taken by the compatibility condition and one new
interface form.

\textbf{Linearisation.} The background $\zeta^0 = \Phi_0\circ\varphi_e$
stays single-valued ($\varphi_e$ is globally smooth); only the
perturbation breaks. With $\varphi^1 := \Phi_1\circ\varphi_e$
(single-valued) and $\mathbf{b} := F_e^{-T}\nabla\zeta^0 =
\nabla\Phi_0\circ\varphi_e$ (continuous across $\Sigma$, since
$\Phi_0$ is $C^1$) as in \S\ref{sec:gravhess},
\begin{equation}
  \zeta_r^1 = \varphi^1 + \mathbf{b}\cdot\bv_r
  \qquad\Longrightarrow\qquad
  \boxed{\ \jump{\zeta^1} = \mathbf{b}\cdot\jump{\bv}
        = -\,\bs\cdot\nabla_\Sigma\zeta^0 \ }
  \quad\text{on each } \Sigma,
  \label{eq:zjump}
\end{equation}
using $\jump{\bv} = -F_e\bs$ and $\mathbf{b}\cdot F_e\bs =
\nabla\zeta^0\cdot\bs$, tangential because $\bs$ is. The first form
$\jump{\zeta^1} = \mathbf{b}\cdot\jump{\bv}$ is the one to impose: a
homogeneous linear constraint pairing the scalar-trace jump with the
displacement-trace jump, needing neither tangential projection nor
the normal-jump constraint --- it joins $\bnu\cdot\jump{\bv} = 0$ in
the same penalty + augmented-Lagrangian machinery, through the same
trace pairing $J = \Pi_{\mathrm{s}}^T\Pi_{\mathrm{f}}$.

\textbf{Second variation.} Differentiating \eqref{eq:Ebroken} twice
along the constrained families of \S\ref{sec:kin}, two new interface
items appear relative to the single-valued route, and only these ---
the sources are linear in $\zeta$ and carry no motion. (i) The
second-order compatibility forces a jump on $\zeta^2$,
\begin{equation}
  \jump{\zeta^2}
  = -\bigl( \nabla\zeta^0\cdot\bm{\sigma}_2
          + \nabla\nabla\zeta^0[\bs,\bs]
          + 2\,\nabla_\Sigma\zeta_{\mathrm{s}}^1[\bs] \bigr),
\end{equation}
whose energy pairing, after the per-region background equation kills
the interior (the envelope structure of $\mathbf{u}_2$ again: free
interior values are gauge, the forced jump is physics), is the flux
term $\oint_\Sigma q\,(-\jump{\zeta^2}) \dS$ with the continuous
background flux density
\begin{equation}
  q := \frac{1}{4\pi G}\,(a_e\nabla\zeta^0)\cdot\bN
     = \frac{1}{4\pi G}\,\mathbf{b}\cdot\bnu .
\end{equation}
(ii) The fluid's second-order boundary datum $\mathbf{b}_2$ of
\eqref{eq:expand} now pairs not only with the elastic traction
$P_{\mathrm{f}}\bN = -\pi\bnu$ (giving $B_\Sigma$ exactly as before)
but also with the gravitational Piola--Maxwell stress of the
$a$-form,
\begin{equation}
  T_g := \frac{J_e}{8\pi G}\Bigl[\,|\mathbf{b}|^2 F_e^{-T}
        - 2\,\mathbf{b}\otimes(F_e^{-1}\mathbf{b})\Bigr],
  \qquad
  \Div T_g = -\rho\,\mathbf{b},
  \qquad
  T_g\bN = \frac{|\mathbf{b}|^2}{8\pi G}\,\bnu - q\,\mathbf{b},
\end{equation}
since in this organisation the gravity--motion coupling sits in the
field term: $\tfrac{1}{8\pi G}\langle a'(\bm{\eta})\nabla\zeta^0,
\nabla\zeta^0\rangle = T_g : D\bm{\eta}$, and the divergence identity
is what re-assembles the interior equilibrium ($\Div(P + T_g) = 0$)
so that the free interior of $\mathbf{u}_2$ still drops. In the sum
of (i) and (ii) the parameterisation freedom cancels twice over: the
\emph{entire} $\bm{\sigma}_2$-dependence --- free tangential part
$\bs_2$ and forced curvature part alike --- cancels between the
$\nabla\zeta^0\cdot\bm{\sigma}_2$ of the jump datum and the
$\mathbf{b}\cdot F_e\bm{\sigma}_2$ of the Maxwell traction, and the
remaining second derivatives collapse by the mechanism of
Remark~\ref{rem:wein} together with the change of variables
\eqref{eq:zjump}. (The first variation performs the same
cancellation one order down: $\oint (T_g\bN)\cdot F_e\bs$ against the
first-order jump flux $\oint q\,\bs\cdot\nabla_\Sigma\zeta^0$ ---
equilibrium is preserved in the broken organisation, as it must be.)

\begin{proposition}[The gravity interface form]\label{prop:GSigma}
In the broken-$\zeta$ organisation the entire gravitational interface
contribution to the second variation is, in Hessian normalisation,
\begin{equation}
  G_\Sigma
  = \oint_\Sigma \Bigl[\,
      \frac{|\mathbf{b}|^2}{8\pi G}\,
        \bnu\cdot\nabla_\Sigma(\bv_{\mathrm{s}} + \bv_{\mathrm{f}})[\bs]
      + q\,\Bigl( \nabla_\Sigma(\zeta_{\mathrm{s}}^1
                          + \zeta_{\mathrm{f}}^1)[\bs]
      - \mathbf{b}\cdot\nabla_\Sigma(\bv_{\mathrm{s}}
                          + \bv_{\mathrm{f}})[\bs] \Bigr)
    \Bigr] \dS,
  \qquad \bs = -F_e^{-1}\jump{\bv},
  \label{eq:GSigma}
\end{equation}
used, like \eqref{eq:BSigmaPol}, in its symmetrised polarisation over
the product space (contributing $(\bv,\bv')$-, $(\zeta^1,\bv')$- and
$(\bv,\zeta^{1\prime})$-blocks). Coefficients: $\bnu$, $\mathbf{b}$,
$q$ --- background data and \emph{tangential} derivatives of the four
traces only; no curvature, no $D^2\varphi_e$, no $\nabla\nabla\zeta^0$.
The vector--vector part is $B_\Sigma$-shaped (mean trace-derivative
against the slip) with $-\pi\bnu$ replaced by
$\tfrac{|\mathbf{b}|^2}{8\pi G}\bnu - q\,\mathbf{b}$; the
scalar--vector part, the mean scalar trace-derivative against the
slip, is the interface footprint of the vanished mismatch coupling
\eqref{eq:mismatch}.
\end{proposition}

Two derivations of \eqref{eq:GSigma} agree: the referential collapse
sketched above, and substituting $\zeta_r^1 = \varphi^1 +
\mathbf{b}\cdot\bv_r$ to reach the Eulerian intermediate form
$\oint q\,[\,2\nabla_\Sigma\varphi^1[\bs] +
D\mathbf{b}[\bs]\cdot(\bv_{\mathrm{s}}+\bv_{\mathrm{f}})\,] +
\oint \tfrac{|\mathbf{b}|^2}{8\pi G}\bnu\cdot\nabla_\Sigma(\bv_{\mathrm{s}}
+ \bv_{\mathrm{f}})[\bs]$ and converting back.

\textbf{The collected broken operator.} Unknowns
$(\bv_{\mathrm{s}}, \bv_{\mathrm{f}}, \zeta_{\mathrm{s}}^1,
\zeta_{\mathrm{f}}^1)$ (the outer scalar field living on solid +
buffer). Volume: the referential operators \emph{per region with each
region's own field} --- elastic/geometric blocks, the mapped Poisson
block per region with the DtN on the outer one, and the
$a'$-coupling $c_r(\zeta_r^1, \bv_r)$ --- no mismatch terms, no
$\rho\nabla\nabla\zeta^0$ mass term, no extension terms off the
buffer. Interface, \emph{per interface}: $B_\Sigma$ unchanged,
$G_\Sigma$ of \eqref{eq:GSigma}, and the two trace constraints
$\bnu\cdot\jump{\bv} = 0$ and $\jump{\zeta^1} =
\mathbf{b}\cdot\jump{\bv}$, both through the same one-sided pairing.
Limits: welded ($\bs \equiv 0$) gives $G_\Sigma = 0$ and
$\jump{\zeta^1} = 0$ --- the continuous-$\zeta$ welded problem
exactly; gravity off ($\mathbf{b} = 0$) gives $G_\Sigma = 0$,
$\jump{\zeta^1} = 0$ and the potential decouples --- the sliding
tests. Topology: everything above is per-interface and no
interpolation between independent solid motions ever arises --- the
trivial scaling promised. The organisation is exactly equivalent to
the single-valued one (bijective change of variables in the
stationary problem), and a solve run both ways is a genuinely
two-sided check --- the two share \emph{no} gravity-interface
machinery: mismatch volume terms + extension on one side, $G_\Sigma$
+ jump condition on the other.

\section{The collected operator}
\label{sec:collect}

Unknowns $(\bv_{\mathrm{s}}, \bv_{\mathrm{f}}, \zeta^1)$ with the
constraint \eqref{eq:constraint}, enforced as in \S\ref{sec:enforce};
this is the single-valued organisation, the broken-$\zeta$ alternative
being collected at the end of \S\ref{sec:brokenzeta}:
\begin{itemize}
\item volume: the referential operators per region with each region's
  own field (fluid dictionary by Lemma~\ref{lem:mub}), the mapped
  Poisson $\zeta$-block and DtN untouched (single-valued $\zeta$,
  $\varphi_e = \mathrm{id}$ at the sphere);
\item the fluid-mismatch gravity coupling of Lemma~\ref{lem:vol},
  assembled with the extension field (extension = gauge; invariance
  testable);
\item the single interface bilinear form $B_\Sigma$ of
  Proposition~\ref{prop:BSigma}: coefficients $\pi$, $\bnu$, $F_e$ and
  \emph{tangential} derivatives of the two traces only (surface
  shape-function derivatives suffice; no $D^2\varphi_e$, no curvature),
  assembled once on the solid side through the pairing of
  \S\ref{sec:disc}, like the constraint itself.
\end{itemize}

Four reductions follow from the construction, and each is checked at
solution level in \S\ref{sec:verif} (in addition to the
finite-difference verification of every coefficient against the exact
broken-motion energy, for which the constrained parameterisation of
\S\ref{sec:kin} satisfies the attachment exactly at finite $\eps$):
(a) spherical hydrostatic reference $\to$ the F1--F3 interface terms
of \texttt{doc/self\_gravitation.md} under the change of variables
$\zeta^1 = \varphi^1 + \bv\cdot\nabla\Phi_0$, and the forms of AC18
(eqs.~121/137) \cite{AC18} and Valette \cite{Valette1986,Valette1991};
(b) welded limit $\bs = 0$: $B_\Sigma = 0$, $\bw$ interior-gauge --- the
welded referential problem of \texttt{doc/gravitating\_elasticity.md}
exactly; (c) gravity and pressure off: the gravity-free sliding
interface; (d) the broken-$\zeta$ organisation of
\S\ref{sec:brokenzeta} against the single-valued solve, head to head
--- the mismatch self-benchmark.

\section{Discretisation of the slipping interface}
\label{sec:disc}

\subsection{The broken space}

The discrete displacement is sought in the \emph{broken} space
\begin{equation}
  V_{\mathrm{s}} \times V_{\mathrm{f}}, \qquad
  V_{\mathrm{s}} \subset H^{1}(\Bs)^{d}, \quad
  V_{\mathrm{f}} \subset H^{1}(\Bf)^{d},
\end{equation}
two independent $H^{1}$ spaces on the two subdomains --- the
discontinuity costs nothing to represent because nothing couples the
spaces yet. The physics then requires exactly one interface condition on
the displacement, the normal constraint \eqref{eq:constraint}, with the
tangential jump free (the slip). It is a \emph{constraint}, enforced
weakly (\S\ref{sec:enforce}); the tangential freedom needs no action at
all. In contrast with a DG discretisation there is no interior-penalty
stabilisation, no mesh-skeleton terms and no consistency issue: inside
each subdomain the space is conforming, and the only broken surface is
$\Sigma$ itself. In the broken-$\zeta$ organisation the potential is
broken the same way, $(\zeta_o, \zeta_{\mathrm{f}})$ on the outer (solid
+ buffer) and fluid regions, with the scalar-jump constraint
\eqref{eq:zjump}.

\subsection{The nodal pairing \texorpdfstring{$J$}{J}}

Enforcing \eqref{eq:constraint} weakly needs the two interface traces in
one bilinear form, i.e.\ a dictionary between fluid-side and solid-side
interface dofs. Both $\Bs$ and $\Bf$ are SubMeshes of one parent mesh,
and each space carries a signed dof injection into a parent-mesh space
(\texttt{SubMeshDofInjection}),
\begin{equation}
  \Pi_{\mathrm{s}} : V_{\mathrm{s}} \to V_{\mathrm{parent}}, \qquad
  \Pi_{\mathrm{f}} : V_{\mathrm{f}} \to V_{\mathrm{parent}},
\end{equation}
sparse Boolean-up-to-sign matrices (the signs carry the orientation
bookkeeping of vector-valued nodal bases). A dof of $V_{\mathrm{s}}$ and
a dof of $V_{\mathrm{f}}$ describe \emph{the same interface node} exactly
when their parent images coincide, so the product
\begin{equation}
  J \;=\; \Pi_{\mathrm{s}}^{T}\, \Pi_{\mathrm{f}}
  \qquad (\text{solid dofs} \times \text{fluid dofs, entries } \pm 1)
  \label{eq:pairing}
\end{equation}
is the complete identification: it is zero except on interface dofs,
where it pairs each fluid trace dof with its solid counterpart,
orientation included. Nothing geometric is computed --- no point
matching, no tolerance --- because the identification is inherited from
the parent's connectivity (\texttt{NewSubMeshPairingMatrix}). In
parallel the same formula runs on \emph{true} dofs as a hypre sparse
product (\texttt{NewSubMeshPairingTrueDofMatrix}); an interface node
whose solid-side and fluid-side owners are different MPI ranks is
handled by the product's communication with no special casing
(\texttt{doc/submesh\_coupling.md}, section ``Pairing two sibling
submeshes''). The broken-$\zeta$ organisation uses the scalar pairing
$J_\zeta$ built the same way.

Because the two trace spaces sit on the same parent faces with the same
nodal basis, they are the \emph{same} discrete trace space in different
coordinates, and $J$ is (up to sign) a permutation between them. That is
what lets every interface bilinear form be assembled \textbf{once}, on
one side only.

\subsection{Why the coupling is weak, not nodal}

Given \eqref{eq:pairing} one is tempted to impose \eqref{eq:constraint}
\emph{strongly}: constrain the normal component of each paired dof. This
fails at the first corner: on a faceted (or isoparametric) interface the
discrete normal is single-valued on element faces but \emph{ill-defined
at the nodes} where faces meet --- exactly where the dofs sit. Any nodal
choice of normal is arbitrary at the vertices, and the resulting
constraint is inconsistent by $O(h)$. The weak form avoids the issue,
because it only evaluates the normal at quadrature points in face
interiors, where it is unambiguous. With $B_n$ the normal--normal
boundary form on the solid side, built with the Nanson normal
\eqref{eq:nanson} (\texttt{BoundaryNormalNormalIntegrator} with the
equilibrium mapping, which applies \texttt{Diffeomorphism::MapNormal}),
\begin{equation}
  (B_n)_{ij} = \oint_{\Sigma} (\bm{\phi}_{i}\cdot\bnu)
               (\bm{\phi}_{j}\cdot\bnu) \dS ,
\end{equation}
the penalty on the normal jump of the paired field expands, via
$\bv_{\mathrm{f}}|_{\Sigma} = $ ($J$-image of the fluid dofs), into the
$2\times 2$ block matrix
\begin{equation}
  \theta\, P, \qquad
  P \;=\;
  \begin{pmatrix}
    B_n & -B_n J \\
    -J^{T} B_n & \;J^{T} B_n J
  \end{pmatrix},
  \qquad
  \theta\, U^{T} P\, U
  \;=\; \theta \oint_{\Sigma} \bigl( \bnu\cdot\jump{\bv} \bigr)^{2} \dS .
  \label{eq:penaltyblock}
\end{equation}
$P$ is symmetric positive semi-definite by construction; its kernel
contains every field with continuous normal trace --- and every purely
tangential jump, which is the point.

\subsection{One-sided assembly of the interface forms}

The interface forms are assembled the same way. With $G$ the one-sided
kernel of \texttt{SlipInterfacePressureIntegrator} on the solid
boundary elements, $S = [I, J]$ the sum map and $D = [I, -J]$ the jump
map, the polarised form \eqref{eq:BSigmaPol} is
\begin{equation}
  B_\Sigma = -\tfrac12\,(S^T G D + D^T G^T S),
\end{equation}
the minus because $G$ is assembled with the solid SubMesh's outward
boundary normal, which is $-\bN$ of the derivation
(\texttt{NewSlipInterfaceMatrix}). The broken-$\zeta$ form
\eqref{eq:GSigma} combines the vector kernel of
\texttt{SlipInterfaceGravityIntegrator} and the scalar--vector kernel of
\texttt{SlipInterfaceGravityScalarIntegrator} in the same way, with a
plus sign (both $\bnu$-linear coefficients flip with the normal,
cancelling the minus of $\bs = -F_e^{-1}\jump{\bv}$), and the scalar
pairing on the potential slots (\texttt{NewSlipGravityInterfaceMatrix}).
Every builder has a parallel overload forming the blocks by hypre
products with the true-dof pairings.

\subsection{The enlarged null space}

A frictionless axisymmetric interface transmits no torque: with the
discontinuity active, the shell and the core rotate
\emph{independently}, so the rigid null space is larger than the welded
problem's --- common translations $(\mathbf{t}, \mathbf{t}, 0)$ and the
independent (mapped) rotations of shell and core, plus the potential
constant in 2-D --- and all such modes are projected. (Translations of
either component alone are \emph{not} null: gravity couples them.) The
constraints vanish on these pairs, since $\bnu$ and $\mathbf{b}$ are
physically radial against tangential rotations. With the tapered fluid
extension of the single-valued organisation the translations are
near-null rather than exact, as for the vacuum extension. On an
interface that is not axisymmetric (an ellipse) the relative rotation is
a stiff physical mode, because the interface transmits torque through
the normal forces alone (\texttt{examples/sliding\_fluid\_ellipse.cpp}).

\section{Constraint enforcement}
\label{sec:enforce}

The interface constraint $\bnu\cdot\jump{\bv} = 0$ (in the broken
organisation also $\jump{\zeta^1} = \mathbf{b}\cdot\jump{\bv}$) is a
genuine, full-rank condition on the interface traces. Two enforcement
paths are implemented on the same assembly: penalty plus augmented
Lagrangian (AL), the default and the only path of the broken
organisation; and a monolithic KKT (saddle-point) solve with an
explicit multiplier, for the single-valued organisation. The fluid
displacement is in addition determined only up to linearised
relabellings; that indeterminacy is an underdetermination, not an
interface condition, and in \emph{both} paths it is handled by the
gauge penalty $\eps Q$ (\texttt{SetFluidGauge}, the covariant
deviatoric penalty on the fluid space) with iterated Tikhonov
refinement, as in \texttt{doc/gauge\_penalty\_iteration.tex}.

\paragraph{The discrete constrained problem.} Collect the single-valued
problem as follows: $U = (U_{\mathrm{s}}, U_{\mathrm{f}}, Z)$, the
coefficient vectors of $\bv_{\mathrm{s}}$, $\bv_{\mathrm{f}}$ and
$\zeta^1$; $A$ the symmetric physical block operator of
\S\ref{sec:collect} (elastic, geometric, gravity, mismatch and
$B_\Sigma$ blocks); $f$ the load; $\eps Q$ the fluid-gauge penalty on
the $U_{\mathrm{f}}$ block; and $P = D^{T}B_nD$ the constraint penalty
\eqref{eq:penaltyblock} on the displacement blocks, with $D = [I, -J]$
the jump map of \S\ref{sec:disc}. The physical problem is: find $U$ and
a multiplier (a dual vector) $w^*\in\mathrm{range}(P)$ with
\begin{equation}
  A\,U + w^* = f, \qquad P\,U = 0 .
  \label{eq:constrained}
\end{equation}
$w^*$ is the discrete interface normal-traction perturbation paired with
$\bnu$ (\S\ref{sec:first}).

\subsection{Penalty and augmented Lagrangian, interleaved with the gauge
refinement}
\label{sec:al}

A pure penalty, $(A + \theta P)U = f$, enforces $PU = 0$ only as
$\theta \to \infty$, where $A + \theta P$ becomes arbitrarily
ill-conditioned --- the dilemma of a single penalised solve
(\texttt{doc/gauge\_penalty\_iteration.tex}, section ``What one
penalised solve gives'') in its constraint-enforcing guise --- and the
same escape applies: iterate rather than stiffen
\cite{Hestenes1969,Powell1969,FortinGlowinski}.

\paragraph{The multiplier iteration.} The augmented-Lagrangian
(Uzawa-type) iteration keeps $\theta$ moderate and updates a running
multiplier $w$:
\begin{equation}
  (A + \theta P)\,U^{(k+1)} = f - w^{(k)}, \qquad
  w^{(k+1)} = w^{(k)} + \theta P\,U^{(k+1)}, \qquad w^{(0)} = 0 .
  \label{eq:alplain}
\end{equation}
At a fixed point the update stalls, so $PU = 0$, and then
$AU + w = f$: the fixed point is \eqref{eq:constrained}, exact at
finite $\theta$ and independent of it.

\paragraph{Contraction rate.} Let $(U^*, w^*)$ solve
\eqref{eq:constrained}; since $PU^* = 0$, $(A + \theta P)U^* = f - w^*$.
Subtracting, the errors $e_U^{(k)} = U^{(k)} - U^*$,
$e_w^{(k)} = w^{(k)} - w^*$ obey
\[
  (A + \theta P)\,e_U^{(k+1)} = -e_w^{(k)}, \qquad
  e_w^{(k+1)} = e_w^{(k)} + \theta P e_U^{(k+1)}
  = \bigl(I - \theta P(A + \theta P)^{-1}\bigr)e_w^{(k)} .
\]
For the analysis factor $P = R^{T}R$ (e.g.\ $R = B_n^{1/2}D$; never
formed), assume $A$ invertible (on the complement of the projected null
pairs of \S\ref{sec:disc}), and write $e_w = R^{T}\mu$ (all multipliers
lie in $\mathrm{range}(P) = \mathrm{range}(R^{T})$). The
Sherman--Morrison--Woodbury identity gives, with the constraint-space
Schur complement $\hat S := RA^{-1}R^{T}$,
\begin{equation}
  R(A + \theta R^{T}R)^{-1}R^{T} = \hat S(I + \theta\hat S)^{-1},
  \label{eq:woodbury}
\end{equation}
hence
\begin{equation}
  e_w^{(k+1)} = R^{T}\bigl[I - \theta\hat S(I + \theta\hat S)^{-1}\bigr]\mu^{(k)}
             = R^{T}(I + \theta\hat S)^{-1}\mu^{(k)} .
\end{equation}
On an eigenvector of $\hat S$ with eigenvalue $\sigma_j$ the multiplier
error is multiplied per sweep by
\begin{equation}
  \frac{1}{1 + \theta\sigma_j} = \frac{\sigma_j^{-1}}{\sigma_j^{-1} + \theta}
  \;\le\; \frac{\norm{A_{\mathrm{S}}}}{\norm{A_{\mathrm{S}}} + \theta},
  \qquad A_{\mathrm{S}} := \hat S^{-1},
\end{equation}
so $A_{\mathrm{S}}$ is the stiffness that the constraint ``sees'', and
the slowest mode is the constraint direction along which $A$ is
stiffest. The constraint violation
$RU^{(k+1)} = Re_U^{(k+1)} = -\hat S(I + \theta\hat S)^{-1}\mu^{(k)}$
contracts at the same rate. Larger $\theta$ contracts faster per sweep
but worsens the conditioning of $A + \theta P$, i.e.\ the cost of each
sweep: moderate $\theta$ and a few sweeps replace a stiff penalty, on
the same amortised-solver economics as the gauge refinement (one
operator, one preconditioner, several solves).

\paragraph{Interleaving with the gauge refinement.} The multiplier
update and the fluid-gauge refinement run in a single loop:
\begin{equation}
  \bigl( A + \eps Q + \theta P \bigr)\, U^{(k+1)}
     = f \;-\; w^{(k)} \;+\; \eps\, Q\, U^{(k)} ,
  \qquad
  w^{(k+1)} = w^{(k)} + \theta\, P\, U^{(k+1)} ,
  \label{eq:alupdate}
\end{equation}
starting from $w^{(0)} = 0$ and $U^{(0)}$ the previous solution, with
$Q$ acting on the fluid block only; each sweep is one block-MINRES solve,
warm-started from the current iterate. The added $\eps QU^{(k)}$ is the
Tikhonov correction. At a fixed point of \eqref{eq:alupdate} the
multiplier update stalls, which forces $P U = 0$ --- the normal jump
vanishes \emph{exactly} (in the discrete weak sense) at finite
$\theta$ --- and the gauge term cancels, so $U$ solves
\eqref{eq:constrained} with neither $\theta$ nor $\eps$ contaminating
the observables. The running multiplier $w$ is the discrete constraint
reaction, i.e.\ the first-order normal traction paired with $\bnu$
(\S\ref{sec:first}; \texttt{ConstraintMultiplier()}). Each part
contracts at its own rate in isolation; the combined loop contracts the
normal jump by about two orders of magnitude over the default sweeps,
to a floor set by the interleaved gauge source and the solver
tolerance, $\theta$-independent relative to the trace, while the
tangential jump remains $O(1)$: the slip the welded space cannot
represent.

The penalty $\theta$ also conditions the solve: it regularises the
otherwise near-null relative normal motion at the interface, so too
small a $\theta$ degrades the solves as well. Below $\theta \approx 100$
the block solves degrade markedly, above it the cost is flat
(\texttt{doc/benchmarks.tex}, ``Solver settings: measured behaviour'');
$\theta = 100$ with eight sweeps is the default
(\texttt{SetConstraint(theta, sweeps)}), a compromise between the
contraction per sweep and the conditioning of each sweep.

\paragraph{Algorithm (AL path, \texttt{SetConstraint}$(\theta, K)$).}
\begin{enumerate}
\item Assemble $A$, $\eps Q$, $P$; build the solver blocks
  $A + \eps Q + \theta P$ and the block-diagonal preconditioner (the
  problem's own preconditioners rebuilt on the penalised blocks:
  algebraic multigrid with systems options on the fluid block in
  parallel, Gauss--Seidel serially; the shifted mapped Laplacian on
  $Z$); build the projector of the null pairs (\S\ref{sec:disc}).
\item Set $w\leftarrow0$ and take $U$ from the previous solution (warm
  start; \texttt{ResetSolution()} zeroes it).
\item For $k = 0,\dots,K-1$ ($K$ the sweep count, default 8):
  \begin{enumerate}
  \item form the right-hand side $f - w + \eps QU$;
  \item solve with projected block MINRES, starting from the current
    $U$, to the sweep tolerance of \S\ref{sec:sweeps}, giving the new
    $U$;
  \item update $w\leftarrow w + \theta PU$;
  \item record the normal-jump energy $(j^{T}B_nj)^{1/2}$,
    $j = U_{\mathrm{s}} - JU_{\mathrm{f}}$ (\texttt{NormalJumpHistory()}).
  \end{enumerate}
\item Return $U$ and $w$ (\texttt{ConstraintMultiplier()}).
\end{enumerate}

\paragraph{The broken-$\zeta$ organisation.} The scalar-jump constraint
$c := \jump{\zeta^1} - \mathbf{b}\cdot\jump{\bv} = 0$ joins the loop
with its own penalty $\theta_\zeta P_\zeta$,
\[
  U^{T}P_\zeta U = \oint_\Sigma c^2 \dS ,
\]
assembled from the interface scalar mass $M_z$, the mixed kernel
$K_{vz} = \oint_\Sigma(\mathbf{b}\cdot\bv)\,\zeta\dS$ and the
$\mathbf{b}\otimes\mathbf{b}$ vector mass $P_b$, all on the solid side
through the pairings $J$ and $J_\zeta$, and its own multiplier, updated
in step 3(c) alongside $w$ and sharing the sweep count. The four-block
system $(U_{\mathrm{s}}, U_{\mathrm{f}}, Z_o, Z_{\mathrm{f}})$ is solved
by projected MINRES in the same way, and the per-sweep scalar-jump
energies are recorded (\texttt{ZetaJumpHistory()}).

\subsection{KKT enforcement}
\label{sec:kkt}

\paragraph{The system.} \texttt{EnableKKT()} makes the multiplier an
unknown: $\lambda$, a coefficient vector on the interface boundary dofs
of a scalar space on the solid SubMesh. With $N$ the one-sided flux
kernel $\oint_\Sigma \lambda\,(\bnu\cdot\bv)\dS$ on the solid side, the
constraint row is
\[
  C\,U = N^{T}\bigl(U_{\mathrm{s}} - J\,U_{\mathrm{f}}\bigr)
\]
restricted to those dofs (zero on the potential block). With
$\hat M = \mathrm{diag}(\hat m_i)$ the diagonal of the interface mass
matrix of the multiplier space, one MINRES per outer step solves
\begin{equation}
  \begin{pmatrix} S_\theta & C^{T} \\ C & 0 \end{pmatrix}
  \begin{pmatrix} U \\ \lambda \end{pmatrix}
  =
  \begin{pmatrix} b \\ 0 \end{pmatrix},
  \qquad
  S_\theta := A_\eps + \theta\, C^{T}\hat M^{-1}C,
  \quad A_\eps := A + \eps Q ,
  \label{eq:kkt}
\end{equation}
with $b$ the gauge-refinement right-hand side (below). ($S_\theta$, the
augmented operator block, is not the sum map $S$ of \S\ref{sec:disc}.)
The outer loop then consists of the fluid-gauge refinements alone, in
the fixed-point form of \eqref{eq:alupdate} without the multiplier
update; \texttt{SetConstraint()}'s iteration count bounds it.

\paragraph{The constraint row is covariant by construction.} $N$ is
assembled by \texttt{BoundaryNormalScalarIntegrator} with the
equilibrium mapping and expanded over $(\bv_{\mathrm{s}},
\bv_{\mathrm{f}})$ with $J$. Nanson's relation \eqref{eq:nanson} makes
flux-type boundary forms exact under mappings --- every area factor
cancels --- so the constraint row is exactly the pull-back of its
unmapped counterpart.

\paragraph{The augmentation does not change the solution.} Any solution
of \eqref{eq:kkt} has $CU = 0$, so the augmentation term
$\theta C^{T}\hat M^{-1}CU$ vanishes and $(U,\lambda)$ solves the
unaugmented system, for every $\theta$; $\theta$ affects only the
conditioning. The KKT endpoint is therefore independent of $\theta$ to
solver precision, whereas the AL endpoint moves with $\theta$ through
its constraint floor. This makes $\theta$-independence a sharp
regression check on the constraint blocks.

\paragraph{The multiplier is physical.} By \S\ref{sec:first}, $\lambda$
is the interface normal-traction perturbation paired with $\bnu$
(\texttt{KKTMultiplier()}), and a cross-check against the AL path's
accumulated \texttt{ConstraintMultiplier()}.

\paragraph{Enforcement is weak.} $CU = 0$ zeroes the jump against the
multiplier space, so the $L^2$ jump of the solution is the projection
remainder --- the discrete constraint floor, the same level a converged
AL iteration reaches. The KKT jump is not expected below the AL one; the
two paths differ in route, not destination.

\paragraph{The multiplier space.} The space is passed to
\texttt{EnableKKT()}. Equal order with the displacement is the
strictest constraint space (used by the serial cross-check test). One
order lower is the mortar-style choice, as in mortar methods, where the
multiplier space is chosen coarser than the trace space
\cite{Wohlmuth2001}: it improves the inf-sup margin of the trace
pairing and shrinks the multiplier block, and it is cheaper at an
endpoint shift well below the mesh error --- a change of constraint
discretisation. It is the recommended choice in 3-D. (The Love-number
benchmark's \texttt{-kkt-order 0}, its default, means equal order;
\texttt{-kkt-order 1} with linear displacements is the recommended
setting.) It is not robust in every configuration: in the 2-D
\texttt{examples/slipping\_interface.cpp}, a multiplier one order
below quadratic displacements makes the gauge refinements diverge (the
normal jump grows by orders of magnitude per refinement), while equal
order converges. When using a lower-order multiplier, check that the
normal-jump history decreases.

\paragraph{The augmented blocks, and why their kernel must match.}
Keeping a $\theta$-penalty inside $S_\theta$ (augmented-Lagrangian
preconditioning in the sense of Golub \& Greif \cite{GolubGreif2003})
serves two purposes: it regularises the near-null relative normal motion
at the interface --- the same role as in the AL path, and the reason
lowering $\theta$ makes MINRES worse even though the multiplier owns the
constraint --- and it shapes the constraint Schur complement.
Eliminating $U$ from \eqref{eq:kkt} gives the multiplier Schur
complement $\mathcal{S} = C S_\theta^{-1} C^{T}$. Write
$R = \hat M^{-1/2}C$, so that the augmentation is $\theta R^{T}R$, and
now $\hat S = RA_\eps^{-1}R^{T}$. The Woodbury identity
\eqref{eq:woodbury} (with $A_\eps$ in place of $A$) gives
\begin{equation}
  \hat M^{-1/2}\,\mathcal{S}\,\hat M^{-1/2} = R S_\theta^{-1}R^{T}
  = \hat S(I + \theta\hat S)^{-1},
\end{equation}
whose eigenvalues are $\sigma_j/(1 + \theta\sigma_j)\in(0, 1/\theta)$,
$\sigma_j$ those of $\hat S$. Equivalently, the preconditioned Schur
complement $(\theta\hat M^{-1})\,\mathcal{S}$ is similar to
$\theta\hat S(I + \theta\hat S)^{-1}$, with eigenvalues
\begin{equation}
  \frac{\theta\sigma_j}{1 + \theta\sigma_j}\in(0,1),
  \qquad \to 1 \ \text{ as } \theta\sigma_j\to\infty
\end{equation}
(and exactly 1 on constraint directions where $A_\eps$ is singular, the
Golub--Greif setting \cite{GolubGreif2003}). Constraint directions on
which $A_\eps$ is soft ($\sigma_j$ large) are therefore clustered at 1:
the Schur complement clusters at $\hat M/\theta$. Only directions with
$\theta\sigma_j\lesssim1$ spread the spectrum, and raising $\theta$
moves them towards 1, so the iteration count falls as $\theta$ grows
even though the multiplier, not the penalty, enforces the constraint;
larger values than the default $\theta = 100$ (shared with the AL path)
gain little. The argument needs the augmentation to be built from
\emph{the same} $C$ and $\hat M$ as the constraint row and its
preconditioner. Augmenting instead with $\theta P$ (the AL penalty's
kernel: the full $L^2$ trace form, not its projection onto the
multiplier space) leaves a kernel mismatch of the size of the gap
between the equal-order trace space and the multiplier space, the
identity above no longer holds, and the solve costs markedly more
iterations. The consistent augmentation is assembled like the $B_n$
folds, as $N\hat M^{-1}N^{T}$ on the interface dofs (the other columns
of $N$ are zero) expanded over $(\bv_{\mathrm{s}}, \bv_{\mathrm{f}})$
with $J$, by serial sparse products or hypre
\texttt{ParMult}/\texttt{RAP}.

\paragraph{The preconditioner.} Block-diagonal and symmetric positive
definite, as MINRES requires:
\begin{equation}
  \mathcal{M}_{\mathrm{KKT}} = \mathrm{diag}\bigl(\mathcal{M}_{S},\;
     P_\lambda\bigr), \qquad P_\lambda = \theta\, \hat M^{-1}
  \qquad (\text{entries } \theta/\hat m_i),
\end{equation}
with $\mathcal{M}_S$ the problem's own block preconditioners rebuilt on
the consistently augmented blocks of $S_\theta$ (algebraic multigrid
with systems options on the fluid block in parallel, Gauss--Seidel
serially; the shifted mapped Laplacian on $\zeta$), and $P_\lambda$ the
\emph{inverse} of the Schur approximation $\hat M/\theta$. With exact
$\mathcal{M}_S = S_\theta^{-1}$ and an exact Schur inverse the
preconditioned KKT matrix would have only the three eigenvalues $1$ and
$(1\pm\sqrt5)/2$ \cite{MurphyGolubWathen2000}; the approximations spread
these into clusters. The dimensional check matters: interface mass
entries scale like $h^{d-1}$, so applying $\theta\hat M$ in place of
$\theta\hat M^{-1}$ mis-scales the multiplier block by
$\theta^2\hat m^2$ --- orders of magnitude at benchmark resolution ---
and multiplies the iteration count while still converging, the endpoint
moving only at solver tolerance. A preconditioner change must leave the
endpoint unchanged at solver tolerance; that, together with the
iteration count, is the test of any change to this block. The null
pairs are the rigid pairs of \S\ref{sec:disc} extended by a zero
multiplier component; the constraint row annihilates them.

\paragraph{Algorithm (KKT path).}
\begin{enumerate}
\item Assemble $A$, $\eps Q$, the kernel $N$ and $\hat M$; form $C$ and
  the augmented blocks $S_\theta$; build $\mathcal{M}_{\mathrm{KKT}}$
  and the projector.
\item For $k = 0,\dots,K-1$ ($K$ the count set by
  \texttt{SetConstraint}): solve \eqref{eq:kkt} with
  $b = f + \eps Q\,U^{(k)}$ for $(U^{(k+1)}, \lambda^{(k+1)})$ by
  projected MINRES, starting from the current iterate ($U^{(0)}$ and the
  starting iterate taken from the previous solution), to the sweep
  tolerances of \S\ref{sec:sweeps}; record the normal-jump energy as a
  diagnostic. This is the iterated Tikhonov refinement of the gauge in
  its non-incremental form (the added $\eps QU^{(k)}$ cancels the
  penalty at the fixed point), with the constraint exact in every
  solve.
\item Return $U$ and $\lambda$ (\texttt{KKTMultiplier()}).
\end{enumerate}

\subsection{The sweep schedule}
\label{sec:sweeps}

Both paths repeat a block solve several times: the AL multiplier sweeps
and the KKT gauge refinements. Only the last sweep determines the
endpoint to full accuracy: an error $\eta$ left by an inexact solve in
sweep $k$ enters $w$ as $\theta P\eta$ and is then contracted by the
remaining sweeps like any multiplier error. \texttt{SetSweepTolerance}
$(\tau_{\mathrm{loose}})$ exploits this. With $\tau$ the solver's
relative tolerance, $\tau_\ell = \max(\tau_{\mathrm{loose}},\tau)$, $K$
sweeps and $N_0$ the initial residual norm of the first sweep (in the
solver's own norm), sweep $k = 0,\dots,K-1$ runs to
\begin{align}
  k = 0&: \text{ relative tolerance } \tau_\ell; \nonumber\\
  1\le k<K-1&: \text{ absolute tolerance }
    \max\bigl(\tau,\ \tau_\ell\,d^{\,k}\bigr)\,N_0,
  \qquad d = (\tau/\tau_\ell)^{1/(K-2)};
\end{align}
and the last sweep $k = K-1$ always runs to $\tau N_0$ (to relative
tolerance $\tau$ when $K = 1$). The tolerance thus tightens
geometrically from $\tau_\ell$ to $\tau$ across the sweeps, and every
requested sweep runs. The design reasons:
\begin{itemize}
\item an early sweep need only be solved roughly, because the next
  multiplier update perturbs its right-hand side anyway, and the late
  accurate sweeps let the AL contraction wash out the early sweeps'
  multiplier error;
\item the later sweeps use an \emph{absolute} target anchored to $N_0$
  because they are warm-started: a relative target would chase each
  sweep's shrinking initial residual;
\item the schedule is fixed in advance, not adaptive: an inner
  tolerance coupled to the measured jump would deadlock (the measured
  jump cannot fall below the solver-error floor that the loose
  tolerance itself sets, while the multiplier error is amplified by
  $\theta$);
\item there is no plateau-based early exit: it triggers falsely on the
  same measurement noise and truncates the multiplier updates the AL
  contraction needs.
\end{itemize}
The library default $\tau_{\mathrm{loose}} = 0$ runs every sweep at full
tolerance --- the reproducible-endpoint mode, mandatory for
finite-difference (shift) studies and strict solver comparisons. The
inexact schedule shifts the endpoint an order below the mesh error at
benchmark resolution and roughly halves the cost; the Love-number
benchmark uses it by default (\texttt{-sweep-tol 1e-3}) and forces exact
mode on \texttt{-map-shift} runs.

\subsection{Choosing between the two paths}

Both paths are kept and solve the same discrete problem: they agree
within discretisation error (at degree 0 to solver tolerance;
\S\ref{sec:verif}; the Love-number benchmark, \texttt{doc/benchmarks.tex},
section ``The Love-number family''), and endpoint differences between
sweep counts and paths sit below the mesh error. At benchmark
resolution their costs are comparable: the AL path with few sweeps is
the cheapest, and the KKT path with a one-order-lower multiplier and the
inexact schedule is cheaper than the AL path at its default eight
full-tolerance sweeps. The AL path is the default and serves both
organisations; the KKT path is single-valued only (\texttt{EnableKKT}
and \texttt{EnableBrokenZeta} are mutually exclusive) and provides the
multiplier as a solved-for field with a $\theta$-independent endpoint.

\subsection{Why the fluid gauge stays a penalty}
\label{sec:gaugekkt}

The KKT idea applies in principle to the fluid gauge as well: minimise
$\tfrac12 (Q\bv, \bv)$ subject to the full coupled equations
$A\bv = f$, with $Q$ the unit-scale deviatoric gauge form --- no $\eps$
anywhere, the gauge fixed as the minimum-$Q$-energy representative, and
the physics exact. It is available on the mixed class
(\texttt{\seqsplit{LinearQuasiStaticMixedSelfGravitatingProblem}::EnableGaugeKKT},
the doubled saddle $[\mathrm{diag}(Q,0), A;\, A, 0]$ with the clean $A$
in the constraint rows), but it is not competitive. Its iterates reach
mesh-level physics early and satisfy the stationarity condition, but
MINRES does not converge algebraically under block-diagonal
preconditioning (Gauss--Seidel or algebraic multigrid) nor under the
Murphy--Golub--Wathen composite on the multiplier slot
\cite{MurphyGolubWathen2000}, at a cost far above the penalty path
(\texttt{doc/benchmarks.tex}). The structural reason: the
``constraint'' operator is square and only approximately singular, so
saddle-point preconditioning theory --- which assumes a full-rank
rectangular constraint --- does not apply, and the near-kernel
reappears wherever the scheme needs solves with $A$. KKT pays for a
full-rank interface constraint, as above; for selection over a
near-kernel, the penalty with iterated refinement is the appropriate
regularisation, and it is the production path for the gauge
(\texttt{doc/gauge\_penalty\_iteration.tex}).

\section{Implementation}
\label{sec:impl}

\paragraph{The class.}
\texttt{\seqsplit{LinearQuasiStaticReferentialSelfGravitatingSlipProblem}}
(\texttt{referential\_problem.hpp}, \texttt{src/referential\_problem.cpp})
derives from the referential self-gravitating problem. Its constructor
takes the solid displacement space (on a SubMesh of the ball; the base
class's displacement space), the fluid displacement space (on a SubMesh
of the same parent, sharing the solid space's finite element
collection), the scalar potential space on the ball, the referential
rheology (position-based coefficients valid on both regions; a fluid
region carries $\mu_b = \pi$ through the bare conversion of
Lemma~\ref{lem:mub}), the referential density, the interface pressure
$\pi$, the interface boundary attributes on the \emph{solid} SubMesh,
$G$, the DtN degree, and optionally the background $\zeta^0$ (otherwise
solved from the density over solid and fluid). Serial and parallel run
in one class: in parallel every block is a hypre product on true dofs
with the cross-rank pairing and extension operators.

\paragraph{Set-up.}
\begin{itemize}
\item \texttt{SetPrescribedVacuumExtension} (base class): the vacuum
  extension of the outer region through the buffer, required in both
  organisations.
\item \texttt{SetFluidExtension(E)}: the prescribed fluid extension
  $\tilde{\bv} = E\,\bv_{\mathrm{s}}$ of the single-valued organisation,
  normally \texttt{NewRadialFluidExtension} --- interface trace rows copy
  the solid values exactly through the pairing, interior rows take
  $t(r)\,\bv_{\mathrm{s}}(r_c\hat{\bx})$ with $t = (r/r_c)^p$ (default
  $p = 2$). The interior rule is gauge; two extensions giving the same
  observables is the built-in invariance test.
\item \texttt{SetFluidGauge(mu\_gauge, epsilon)}: the fluid gauge
  penalty, assembled covariantly as
  \texttt{ElasticTensorIntegrator(C, map)} with the pulled-back
  deviatoric tensor and the rheology's equilibrium mapping. The
  base-class \texttt{SetGaugedFluid()} is refused (the fluid has its own
  space here).
\item \texttt{SetConstraint(theta, iterations)}: the penalty $\theta$
  (default $100$) and the number of AL sweeps or KKT gauge refinements
  per \texttt{Solve()} (default $8$).
\item \texttt{EnableKKT(fes\_scalar\_solid)}: the KKT path of
  \S\ref{sec:kkt}; single-valued organisation only.
\item \texttt{EnableBrokenZeta(fes\_zeta\_outer, theta\_zeta)}: the
  broken-$\zeta$ organisation, with $(\zeta_o, \zeta_{\mathrm{f}})$ on
  the outer region (a shadow space covering every attribute except the
  fluid) and the fluid; the DtN stays on the ball space and is folded
  through the outer injection; no fluid extension is needed.
\item \texttt{SetSweepTolerance(loose)}: the inexact sweep schedule of
  \S\ref{sec:sweeps}.
\end{itemize}

\paragraph{Usage} (single-valued, KKT):
\begin{verbatim}
slip.SetPrescribedVacuumExtension(fes_buffer, *Evac);
slip.SetFluidExtension(*Ef);          // NewRadialFluidExtension
slip.SetFluidGauge(mu_gauge, 1e-2);
slip.SetConstraint(100.0, 3);         // theta; outer iterations
slip.EnableKKT(&fes_scalar_solid);    // one order below u recommended
slip.SetSweepTolerance(0.0);          // 0 = reproducible endpoint
\end{verbatim}
In the Love-number benchmark: \texttt{-method slip} (single-valued) or
\texttt{-method slip\_broken}, with \texttt{-theta}, \texttt{-al},
\texttt{-sweep-tol}, and for the KKT path \texttt{-kkt -kkt-order 1}.
\texttt{examples/slipping\_interface.cpp} exposes every option of the
class on one problem (a two-layer self-gravitating body under a surface
load, spherical or relabelled onto an aspherical reference mesh), prints
the constraint histories and the surface harmonic analysis including the
leakage into other degrees, compares with the welded and Dahlen
solutions, and shows the fields in GLVis.

\paragraph{Restrictions.}
\begin{itemize}
\item The single-valued organisation assembles its gravity
  \emph{mismatch} pieces (\S\ref{sec:gravhess}) at
  $\varphi_e = \mathrm{id}$ only, and refuses a non-identity equilibrium
  mapping (\texttt{Diffeomorphism::IsIdentity}), pointing to the broken
  organisation; its elastic, $B_\Sigma$ and referential gravity terms are
  fully mapped. The broken organisation is assembled mapped throughout
  and is the route for mapped and aspherical backgrounds.
\item One fluid region inside a solid shell: the radial fluid extension
  is one-sided (no nested shells with two-sided extensions).
\item The KKT path is single-valued only.
\end{itemize}

\paragraph{Observables.} The solid displacement, and in the fluid the
quantities built from $\mathrm{div}(\rho\bv_{\mathrm{f}})$ and the
interface normal displacement; the fluid displacement itself is gauge.
In the single-valued organisation $\zeta^1 = \varphi^1 +
\tilde{\bv}\cdot\nabla\zeta^0$ carries the extension field off the
solid, so it is extension-invariant (an observable) only on the solid
region; comparisons elsewhere must first convert to the Eulerian
$\varphi^1$. The multipliers (\texttt{ConstraintMultiplier()},
\texttt{KKTMultiplier()}) are the first-order interface normal traction
paired with $\bnu$, not $\pi$ (\S\ref{sec:first}).

\paragraph{Diagnostics.} \texttt{NormalJumpHistory()} and
\texttt{ZetaJumpHistory()} (per-sweep constraint energies),
\texttt{SlipRigidPairResiduals()} and \texttt{BlockNullPairResidual()}
(the null pairs under the assembled operator),
\texttt{FluidDisplacement()}, \texttt{OuterPotential()} and
\texttt{FluidPotential()} (broken mode), and, for the change-of-variables
identity of \texttt{doc/mappings.md}, per-block and per-kernel probes of
the broken operator (\texttt{ApplyBrokenBlock()},
\texttt{ApplyBrokenKernel()}).

\section{Verification}
\label{sec:verif}

Each identity below is asserted by a test; the quoted levels are those
measured at the test settings (the assertions are looser bounds).
Energy identities use exact constrained broken-motion families on a
circular referential $\Sigma$ (rotational slip: the attachment holds
exactly at finite $\eps$, and the family's own curvature exercises the
$\bm{\sigma}_2$ correction), with the gravitational part of the exact
energy in closed form by the shell theorem (plus the exact $m = 1$
multipole energy in the broken case). Solution-level checks use the
two-layer hydrostatic disc, where $\pi(r_c) > 0$, so $B_\Sigma$ and the
gravity interface or mismatch pieces act.

\medskip
\noindent
{\small
\begin{tabular}{@{}>{\raggedright\arraybackslash}p{0.38\textwidth}>{\raggedright\arraybackslash}p{0.31\textwidth}>{\raggedright\arraybackslash}p{0.25\textwidth}@{}}
\hline
Identity & Test & Level \\
\hline
$d^2E/d\eps^2 = Q_{\mathrm{vol}} + B_\Sigma$ at $\varphi_e =
\mathrm{id}$ and a curved radial $\varphi_e$ (non-trivial $\bnu$, $J_e$),
diagonal and polarised & \texttt{\seqsplit{TestSlipInterface}}:
\texttt{\seqsplit{SecondVariationIdentity}}, \texttt{\seqsplit{SecondVariationMapped}} &
$2\times10^{-4}$ relative; $B_\Sigma$ resolved $>50\times$ the residual \\
discrete $B_\Sigma$ against its quadrature form &
\texttt{\seqsplit{DiscreteMatrixMatchesQuadrature}} & converges with order; welded
pairs annihilated to round-off \\
single-valued gravity Hessian \eqref{eq:mismatch}--\eqref{eq:wextra}:
radial squeeze, azimuthal, axisymmetric relabelling, mixed families,
vanishing gravity--slip cross terms & \texttt{\seqsplit{GravityHessianIdentity}} &
$\le 3\times10^{-4}$ \\
discrete mismatch pieces against quadrature &
\texttt{\seqsplit{DiscreteGravityPiecesMatchQuadrature}} & order 1:
$\sim4\times10^{-2}$; order 2: $2$--$5\times10^{-4}$ \\
$d^2E/d\eps^2$ in the broken organisation, $G_\Sigma$ included; joint
annihilation of the axisymmetric relabelling &
\texttt{\seqsplit{BrokenZetaGravityIdentity}} & $3\times10^{-4}$ \\
discrete $G_\Sigma$ against quadrature, at the identity and a curved
radial $\varphi_e$ & \texttt{\seqsplit{DiscreteGravityInterfaceMatchesQuadrature}}
& converges with order; welded pairs annihilated to round-off \\
pairing $J$: nodal trace identification, $J^TJ = I$ on paired dofs &
\texttt{\seqsplit{TestSlidingInterface}}: \texttt{\seqsplit{PairingIdentities}} & machine
precision \\
gravity-free sliding cavity against the condensed rank-one reference;
normal jump contracts, tangential jump free &
\texttt{\seqsplit{CavityMatchesCondensedWithFreeSlip}},
\texttt{\seqsplit{TestSlidingInterfacePar}} & $<5\times10^{-3}$; normal jump
$<10^{-4}$ of the trace \\
reduction (b): slipping against welded gauged solution (spherical,
barotropic); two fluid extensions (taper powers 2, 3) &
\texttt{\seqsplit{TestSlipProblem}}: \texttt{\seqsplit{TwoLayerBarotropicCrossCheck}} &
$1.1\times10^{-3}$ ($u_s$), $7\times10^{-4}$ ($\zeta^1$ on the solid),
order 2; extensions agree to $5\times10^{-6}$ \\
reduction (a): Dahlen-class solution (F1--F3) under
$\zeta^1 = \varphi^1 + \bv\cdot\nabla\Phi_0$, degree-2 load &
\texttt{\seqsplit{ReducesToDahlenOnSphere}} & $6\times10^{-3}$ ($u_s$),
$5\times10^{-3}$ (potential), order 2 \\
reduction (c): massless, pressure-free background $\to$ gravity-free
sliding cavity & \texttt{\seqsplit{PressureFreeReducesToSlidingCavity}} &
$2\times10^{-5}$; potential decouples to zero \\
KKT against AL & \texttt{\seqsplit{KKTMatchesAugmentedLagrangian}} & displacement
within $10^{-3}$; KKT jump at the AL floor \\
reduction (d): broken against single-valued organisation &
\texttt{\seqsplit{BrokenZetaHeadToHead}} & $7\times10^{-5}$ ($u_s$),
$5\times10^{-5}$ ($\zeta^1$ on the solid) \\
null pairs under the full operator, both penalties included, at the
identity and under a relabelling & \texttt{\seqsplit{RigidPairsNearNull}} &
near-null \\
mapped broken solver: interface-fixing relabelling
$\xi = f(r)\,\mathrm{Rot}(\alpha(r))\,x$ with a twist across $\Sigma$
(non-radial $F_e$ on the interface), against the identity description &
\texttt{\seqsplit{BrokenZetaRelabelledEquilibrium}} & $5\times10^{-4}$ ($u_s$),
$2\times10^{-4}$ ($\zeta^1$ on the solid), two orders below the map
amplitude \\
serial against parallel (single-valued and broken; 1, 2, 4 ranks;
orders 1, 2) & \texttt{\seqsplit{TestReferentialProblemPar}} & $10^{-5}$ \\
\hline
\end{tabular}}

\medskip
The mapped broken check is the slipping-interface instance of the
discrete change-of-variables identity of \texttt{doc/mappings.md}
(section ``The discrete change-of-variables identity''): every mapped
piece (per-region $a$-forms and Poisson blocks, couplings, $B_\Sigma$,
$G_\Sigma$, $\mathbf{b}$, the jump constraint) is load-bearing in it.
The two organisations share \emph{no} gravity-interface machinery
(mismatch volume terms and extension on one side, $G_\Sigma$ and jump
condition on the other), and reduction (a) shares no interface
machinery with the Dahlen formulation, so these are two-sided checks.
Externally, the \texttt{slip} and \texttt{slip\_broken} methods of the
Love-number benchmark are compared with \texttt{pyslfp} and with the
welded methods (\texttt{doc/benchmarks.tex}), which also records the
measured cost of the AL and KKT paths.

\appendix

\section{Assessment of the published treatments}
\label{sec:verdict}

Checked against the derivation above, at the fluid--solid interface:

\textbf{Al-Attar, Crawford, Valentine \& Trampert (2018)
\cite{AC18} [AC18]: correct.} Their natural-configuration quadratic action (their eq.~121)
carries the interface terms
$-\!\int\varpi^{(0)}\langle\nabla\sigma\otimes\jump{\bm{\varphi}},
D\bm{\varphi}_1\rangle + \tfrac12\!\int\varpi^{(0)}\langle
\nabla\nabla\sigma\,\jump{\bm{\varphi}}, \jump{\bm{\varphi}}\rangle$;
doubling (their $\mathcal{S}^{(2)}$ is one half of the Hessian) and
applying the collapse \eqref{eq:collapse} in reverse, this is
\emph{exactly} \eqref{eq:BSigma} at $\varphi_e = \mathrm{id}$ with
$\varpi^{(0)}$ identified with the referential pressure $\pi$ ---
factor of two, curvature term ($\nabla\nabla\sigma = \mathrm{II}$ for
their distance-type $\sigma$) and all. Their general-configuration
operators $(\mathbf{Q}_1^{(0)}, \mathbf{S}_1^{(0)})$ (their
eqs.~107/109) are the pre-collapse form of the same object, and their
own reduction to Woodhouse \& Dahlen (1978) \cite{WoodhouseDahlen1978}
stands. Since our
derivation is independent (constrained-family route, no level-set
machinery), the agreement is a genuine two-sided check.

\textbf{Maitra \& Al-Attar (2024) \cite{MA24} [MA24]: exact theory
correct;
one identified oversight in the linearisation.} Their nonlinear
treatment carries every Jacobian explicitly: the weak slip term (their
eq.~170) pairs against $J\mathbf{F}^{-T}\hat{\mathbf{n}} = \bnu$, the
multiplier identification (their eq.~174) gives the traction
$-\varpi\,\bnu$, i.e.\ $\varpi = \pi$ in agreement with
\S\ref{sec:first}, and the finite-slip traction continuity (their
eq.~176) even displays the surface Jacobian $J_\chi$ with commentary.
The linearised slip vector (their eq.~B3) and constraint (B4) agree
with \S\ref{sec:kin} (the latter is ours divided by $J^{(0)}$ --- the
same constraint set). The oversight sits in their eqs.~(B5)--(B6):
expanding (170) exactly, the $O(\eps)$ terms inherit the factor
$J^{(0)}$ from $\delta(J\mathbf{F}^{-T})$, but their operator
$\mathbf{Q}\,\mathbf{a} = \mathbf{F}^{(0)-T}[\hat{\mathbf{n}}\otimes
\mathbf{a}]\mathbf{F}^{(0)-T} = (\mathbf{F}^{-T}\hat{\mathbf{n}})
\otimes(\mathbf{F}^{-1}\mathbf{a})$ carries no $J^{(0)}$ while their
$O(1)$ term keeps it: \textbf{the two $\mathbf{Q}$-terms of (B6), and
hence the linearised momentum equation using them, are missing a
factor $J^{(0)}$}. (Their simultaneous dropping of the
$\tr(\mathbf{F}^{-1}D\mathbf{u})$-term is, by contrast, legitimate:
it multiplies $\langle\jump{\mathbf{w}},
\mathbf{F}^{(0)-T}\hat{\mathbf{n}}\rangle$, which vanishes on their
constrained test space (B7a).) The oversight is invisible in the
natural configuration ($J^{(0)} = 1$) --- hence in all their stated
applications --- and consistent with the general rule of
Remark~\ref{rem:jacobian}: an oversight, not a convention.

Organisationally the two papers differ: AC18's final weak form
(their eqs.~118--120) imposes the slip constraint on the configuration
space and \emph{eliminates} the first-order multiplier
$\varpi^{(1)}$, so their linearised equations carry only the
\emph{given} background field $\varpi^{(0)}$ (their footnote relates
$-\varpi^{(0)}$ to Woodhouse--Dahlen's $\pi_0$; the identification
with our $\pi$ is up to that sign convention, which nothing here
relies on). MA24 instead imposes the test-function constraint weakly
with a multiplier field $\varpi$ (their eq.~170) --- a mixed
formulation. Our $B_\Sigma$ is AC18's organisation --- background
data only, constraint on the space (via penalty + AL in practice) ---
plus the collapse to trace-only tangential derivatives, which is what
makes the implementation curvature-free.

\begin{thebibliography}{99}

\bibitem{AC18}
D.~Al-Attar, O.~Crawford, A.~P. Valentine, J.~Trampert,
\emph{Hamilton's principle and normal mode coupling in an aspherical
planet with a fluid core},
Geophys. J. Int. 214 (2018), 485--507.

\bibitem{MA24}
M.~Maitra, D.~Al-Attar,
\emph{On the elastodynamics of rotating planets},
Geophys. J. Int. 237 (2024), 1301--1338.

\bibitem{Valette1986}
B.~Valette,
\emph{About the influence of pre-stress upon adiabatic perturbations of
the Earth},
Geophys. J. R. astr. Soc. 85 (1986), 179--208.

\bibitem{Valette1991}
B.~Valette, on the gravito-elastodynamics of a pre-stressed elastic
Earth, Geophys. J. Int. 104 (1991), 555.

\bibitem{WoodhouseDahlen1978}
J.~H. Woodhouse, F.~A. Dahlen,
\emph{The effect of a general aspherical perturbation on the free
oscillations of the Earth},
Geophys. J. R. astr. Soc. 53 (1978), 335--354.

\bibitem{Hestenes1969}
M.~R. Hestenes,
\emph{Multiplier and gradient methods},
J. Optim. Theory Appl. 4 (1969), 303--320.

\bibitem{Powell1969}
M.~J.~D. Powell,
\emph{A method for nonlinear constraints in minimization problems},
in R.~Fletcher (ed.), \emph{Optimization}, Academic Press, 1969.

\bibitem{FortinGlowinski}
M.~Fortin, R.~Glowinski,
\emph{Augmented Lagrangian Methods},
North-Holland, 1983.

\bibitem{GolubGreif2003}
G.~H. Golub, C.~Greif,
\emph{On solving block-structured indefinite linear systems},
SIAM J. Sci. Comput. 24 (2003), 2076--2092.

\bibitem{MurphyGolubWathen2000}
M.~F. Murphy, G.~H. Golub, A.~J. Wathen,
\emph{A note on preconditioning for indefinite linear systems},
SIAM J. Sci. Comput. 21 (2000), 1969--1972.

\bibitem{Wohlmuth2001}
B.~I. Wohlmuth,
\emph{Discretization Methods and Iterative Solvers Based on Domain
Decomposition},
Springer, 2001.

\end{thebibliography}

\end{document}
